%------------------------------------------------------------------------------%
\begin{question}
  Resolva pelo Método de Cramer
  \[
    \begin{cases}
      x   +y + z = 6 \\
      2x  -y + z = 3 \\
      x  +2y - z = 2
    \end{cases}
  \]
\end{question}

%------------------------------------------------------------------------------%
\begin{resolution}{Solução}
  \onslide<+->{\[
    A =
    \begin{pmatrix}
      1 & 1  & 1  \\
      2 & -1 & 1  \\
      1 & 2  & -1
    \end{pmatrix}
  \]}

  \onslide<+->{Pela Regra de Sarrus}
  \begin{align*}
    \onslide<+->{\det(A) & =}
    \onslide<+->{    1(-1)(-1)}
    \onslide<+->{  + 1(1)(1) }
    \onslide<+->{  + 1(2)(2)}
    \onslide<+->{  - 1(-1)(1)}
    \onslide<+->{  - 1(2)(-1)}
    \onslide<+->{  - 1(1)(2)}
    \\
    \onslide<+->{& = 1 + 1 + 4 + 1 + 2 - 2 \\}
    \onslide<+->{& = 7}
  \end{align*}
  \onslide<+->{Como
  \(
    \det(A) \neq 0
  \)}
  \onslide<+->{o sistema possui uma única solução}
\end{resolution}

%------------------------------------------------------------------------------%
\begin{resolution}{Cálculo do Determinante de $A_x$}
  \onslide<+->{\[
    A_x =
    \begin{pmatrix}
      6 & 1  & 1  \\
      3 & -1 & 1  \\
      2 & 2  & -1
    \end{pmatrix}
  \]}
  \begin{align*}
    \onslide<+->{\det(A_x)}
    \onslide<+->{& = 6(-1)^{1+1}
      \begin{vmatrix}
        -1 & 1  \\
        2  & -1
      \end{vmatrix}}
    \onslide<+->{  + 1 (-1)^{1+2}
      \begin{vmatrix}
        3 & 1  \\
        2 & -1
      \end{vmatrix}}
    \onslide<+->{  + 1 (-1)^{1+3}
      \begin{vmatrix}
        3 & -1 \\
        2 & 2
      \end{vmatrix}}
    \\
    \onslide<+->{& = 6\big[(-1)(-1)  -1 \times 2\big]}
    \onslide<+->{  -  \big[3(-1) - 1 \times 2\big]}
    \onslide<+->{  +  \big[3\times 2 - (-1)2 \big] \\}
    \onslide<+->{& = -6 + 5 + 8  \\}
    \onslide<+->{& = 7}
  \end{align*}
\end{resolution}

%------------------------------------------------------------------------------%
\begin{resolution}{Cálculo do Determinante de $A_y$}
  \onslide<+->{\[
    A_y =
    \begin{pmatrix}
      1 & 6 & 1  \\
      2 & 3 & 1  \\
      1 & 2 & -1
    \end{pmatrix}
  \]}
  \begin{align*}
    \onslide<+->{\det(A_y)}
    \onslide<+->{& = 1 (-1)^{1+1}
        \begin{vmatrix}
          3 & 1  \\
          2 & -1
        \end{vmatrix}}
    \onslide<+->{  + 6 (-1)^{1+2}
        \begin{vmatrix}
          2 & 1  \\
          1 & -1
        \end{vmatrix}}
    \onslide<+->{  + 1 (-1)^{1+3}
        \begin{vmatrix}
        2&3\\
        1&2
        \end{vmatrix} \\}
    \onslide<+->{& = 1\big[3(-1) - 1 \times 2\big] }
    \onslide<+->{  - 6\big[2(-1) - 1 \times 1\big]}
    \onslide<+->{  +  \big[2 \times 2 - 3 \times 1\big] \\}
    \onslide<+->{& = -5+18+1  \\}
    \onslide<+->{& = 14}
  \end{align*}
\end{resolution}

%------------------------------------------------------------------------------%
\begin{resolution}{Cálculo do Determinante de $A_z$}
  \onslide<+->{\[
    A_z =
    \begin{pmatrix}
      1 & 1  & 6 \\
      2 & -1 & 3 \\
      1 & 2  & 2
    \end{pmatrix}
  \]}
  \begin{align*}
    \onslide<+->{\det(A_z)}
    \onslide<+->{& = 1 (-1)^{1+1}
        \begin{vmatrix}
          -1 & 3 \\
          2  & 2
        \end{vmatrix}}
    \onslide<+->{  + 1 (-1)^{1+2}
        \begin{vmatrix}
          2 & 3 \\
          1 & 2
        \end{vmatrix}}
    \onslide<+->{  + 6 (-1)^{1+3}
        \begin{vmatrix}
          2 & -1 \\
          1 & 2
        \end{vmatrix} \\}
    \onslide<+->{& = 1\big[(-1)2 - 3 \times 2\big] }
    \onslide<+->{  -  \big[2 \times 2 - 3 \times 1\big]}
    \onslide<+->{  + 6\big[2 \times 2 - (-1)1\big] \\}
    \onslide<+->{& = -8 -1 +30 \\}
    \onslide<+->{& = 21}
  \end{align*}
\end{resolution}

%------------------------------------------------------------------------------%
\begin{resolution}{Solução}
\[
  x
  \pause
  \;=\; \frac{\det(A_x)}{\det(A)}
  \pause
  \;=\; \frac{7}{7}
  \pause
  \;=\; 1
\]

\pause
\[
  y
  \pause
  \;=\; \frac{\det(A_y)}{\det(A)}
  \pause
  \;=\; \frac{14}{7}
  \pause
  \;=\; 2
\]

\pause
\[
  z
  \pause
  \;=\; \frac{\det(A_z)}{\det(A)}
  \pause
  \;=\; \frac{21}{7}
  \pause
  \;=\; 3
\]
\end{resolution}

%------------------------------------------------------------------------------%
\endinput
